\documentclass{beamer} 
\usepackage{beamerthemesplit} 
\usepackage{wrapfig} 
\usepackage{verbatim} 
\usetheme{SPbGU} 
\usepackage{pdfpages} 
\usepackage{amsmath} 
\usepackage{cmap}
\usepackage{array} 
\usepackage[T2A]{fontenc} 
\usepackage[utf8]{inputenc} 
\usepackage{tikz} 
\usetikzlibrary{positioning, automata}
\usepackage{multirow} 
\usepackage[noend]{algpseudocode} 
\usepackage{algorithm} 
\usepackage{algorithmicx} 
\usetikzlibrary{shapes,arrows} 
\usepackage{fancyvrb} 
\usepackage{pgfplots} 
\usepackage{sidecap} 
\usepackage{fontawesome}
\usepackage{soul}
\usepackage{xcolor}
\usepackage{tabu}
\usepackage{tikz}
\usetikzlibrary{calc}
\usepackage{zref-savepos}
\usepackage{colortbl}
\usepackage[normalem]{ulem}
\pgfplotsset{compat=1.9} 
\newtheorem{rutheorem}{Теорема} 
\newtheorem{ruproof}{Доказательство} 
\newtheorem{rudefinition}{Определение} 
\newtheorem{rulemma}{Лемма} 
\beamertemplatenavigationsymbolsempty 

\newcounter{NoTableEntry}
\renewcommand*{\theNoTableEntry}{NTE-\the\value{NoTableEntry}}

\newcommand*{\strike}[2]{%
	\multicolumn{1}{#1}{%
		\stepcounter{NoTableEntry}%
		\vadjust pre{\zsavepos{\theNoTableEntry t}}% top
		\vadjust{\zsavepos{\theNoTableEntry b}}% bottom
		\zsavepos{\theNoTableEntry l}% left
		\hspace{0pt plus 1filll}%
		#2% content
		\hspace{0pt plus 1filll}%
		\zsavepos{\theNoTableEntry r}% right
		\tikz[overlay]{%
			\draw
			let
			\n{llx}={\zposx{\theNoTableEntry l}sp-\zposx{\theNoTableEntry r}sp-\tabcolsep},
			\n{urx}={\tabcolsep},
			\n{lly}={\zposy{\theNoTableEntry b}sp-\zposy{\theNoTableEntry r}sp},
			\n{ury}={\zposy{\theNoTableEntry t}sp-\zposy{\theNoTableEntry r}sp}
			in
			(\n{llx}, \n{lly}) -- (\n{urx}, \n{ury})
			;
		}% 
	}%
}

\title[Complexity of CFL-reachability]{Complexity of context-free language (CFL) reachability problem: new challenges and opportunities} 
% То, что в квадратных скобках, отображается в левом нижнем углу. 
\institute[]{ JetBrains Research, Programming Languages and Tools Lab  \\
    SPBAU, SPBSU } 

\author[Ekaterina Shemetova]{Ekaterina Shemetova}

\date{17.12.2021} 

\begin{document} 
\definecolor{red}{RGB}{255,0,0} 

{
\begin{frame}[fragile]
  \begin{tabular}{p{2.0cm} p{4.5cm} p{1.5cm} p{1cm}}
   \begin{center}
      \includegraphics[height=1.5cm]{pictures/jetbrainsResearch.pdf}
    \end{center}
    &
    \begin{center}
      \includegraphics[height=1.5cm]{pictures/YC_logo.pdf}
    \end{center}
    &
    \begin{center}
      \includegraphics[height=1.5cm]{pictures/au-logo-full.png}
    \end{center}
    &
    \begin{center}
      \includegraphics[height=1.5cm]{pictures/SPbGU_Logo.png}
    \end{center}
  \end{tabular}
  \titlepage
\end{frame}
}


\begin{frame}
\frametitle{CFL-reachability}
        \begin{center}
           graph reachability + constraints in the terms of context-free language
        \end{center}



\tikzset{every state/.style={minimum size=0pt}}
\begin{figure}[h]
    \centering        
    \begin{tikzpicture}[shorten >=1pt,auto, node distance=0.8cm]
       \node[state] (q_0)                      {$s$};
       \node[state] (q_1) [right=of q_0] {};
       \node[state] (q_2) [right=of q_1]       {};
       \node[state] (q_3) [right=of q_2]       {};
       \node[state] (q_4)   [right=of q_3]                    {};
       \node[state] (q_5) [right=of q_4] {};
       \node[state] (q_6) [right=of q_5]       {};
       \node[state] (q_7) [right=of q_6]       {$t$};
        \node[state] (q_8) [above=of q_2]       {};
        \node[state] (q_9) [above right=of q_8]       {};
       \node[state] (q_10) [below right=of q_9]       {};
        \path[->]
        (q_0) edge node {$[$} (q_1)
        (q_1) edge node {$($} (q_2)
        (q_2) edge node {$e$} (q_3)
        (q_3) edge node {$[$} (q_4)
        (q_4) edge node {$e$} (q_5)
        (q_5) edge node {$]$} (q_6)
        (q_6) edge node {$]$} (q_7)
        (q_4) edge[above] node {$]$} (q_10)
        (q_10) edge[above] node {$)$} (q_9)
        (q_9) edge[above] node {$e$} (q_8)
        (q_8) edge[left] node {$e$} (q_2);

    \end{tikzpicture}
\end{figure}
\centering
\textcolor{green}{[(e[])eee[e]]}


\begin{itemize}
\item Static code analysis: interprocedural dataflow analysis, points-to, ..., (equivalence to the subclass of set constraints)
\item Datalog chain queries
\item Graph databases
\end{itemize}
\end{frame}

\begin{frame}
\frametitle{Motivation}
There is no algorithm for the CFL-reachability problem faster than $O(n^3)$ ($O(n^3/\log n$)). 
\begin{block}{Nevin Heintze and David McAllester, 1997}
On the Cubic Bottleneck in Subtyping and Flow Analysis
\end{block}

\begin{itemize}
\item It is an open problem whether truly subcubic $O(n^{3-\varepsilon})$ algorithm exists
\item Under the \textbf{combinatorial BMM-hypothesis} such an algorithm is unlikely to exist...
\item Moreover, CFL-reachability problem is inherently sequential (is hard to parallelize effectively)
\end{itemize}
\end{frame}

\begin{frame}
\frametitle{Plan}
While it is hard to obtain a good solution in the general case, we can go deeper into:
\begin{itemize}
\item Parametrized complexity: fix some parameter (specific language or graph) and try to improve this variant of the problem
\item Fine-grained complexity: no improvement, just better understanding why problem is hard
\item Dynamic complexity: how fast can we update the output when the input is slightly changed
\end{itemize}
\end{frame}



\begin{frame}
\frametitle{Parametrized complexity: languages of bounded dimension}
\begin{block}{Advantages and disadvantages of Linear Datalog}

\faThumbsOUp\ Linear Datalog=Linear languages can be evaluated effectively

\faThumbsOUp\ Linear Datalog=Linear languages is effectively parallelizable

\faThumbsODown\ Linear Datalog=Linear languages is not expressive enough (i.e. for static code analysis)

\end{block}

In order to increase expressivity, we considered the \textbf{languages of bounded dimension} --- generalization of linear languages

\begin{figure}
\centering
\begin{tikzpicture}[
level 1/.style={sibling distance=3cm, level distance=1cm},
level 2/.style={sibling distance=1.5cm}]
%\tikzstyle{every node}=[circle,draw]

\node[circle,draw] (Root) [ fill=red, red] {}
    child {
    node[circle,draw, blue] (l) {} 
    child { node[circle,draw, fill=red, red](ll) {}
            child { node[circle,draw, blue] (p) {} }
            child { node[circle,draw, blue] (pl) {} }
             }
    child { node[circle,draw, blue](lr) {} 
          child { node[circle,draw, blue] (plr) {} }
      }
}
child {
    node[circle,draw,  fill=red, red] (r) {}
    child { node[circle,draw, blue] (rl) {}} 
    child { node[circle,draw, blue] (rr) {} }
};
\node  [right=0.05cm of p] {0};
\node  [right=0.1cm of Root] {$2, dim(t)=2$};
\node  [right=0.05cm of l] {1};
\node  [right=0.05cm of r] {1};
\node  [right=0.05cm of ll] {1};
\node  [right=0.05cm of lr] {0};
\node  [right=0.05cm of pl] {0};
\node  [right=0.05cm of plr] {0};
\node  [right=0.05cm of rl] {0};
\node  [right=0.05cm of rr] {0};
\end{tikzpicture}
\end{figure}

\end{frame}

%-------------------------------------------------------------------------------------------------------------
\begin{frame}
\frametitle{Parametrized complexity: languages of bounded dimension}
\textbf{Our results:}
\begin{itemize}
\item CFL-reachability where input is the language of bounded dimension can be effectively evaluated and parallelized 
\item The rational index (complexity measure of language in FL theory) of the languages of bounded dimension $d$ is no more than $O(n^{2d})$
\item Will be submitted to DLT/SCR/journal
\end{itemize}

\textbf{Future work:}
\begin{itemize}
\item Generalize our techniques to estimate the complexity of the reachability in terms of MCFG (multiple CFG)
\item Languages of bounded dimension can be more precise and still effective approximation of analysis in compare with the linear ones --- estimate what variants of CFL-reachability based static code/data analyses can be effectively approximated in such a way
\end{itemize}
\end{frame}
%-------------------------------------------------------------------------------------------------------------
\begin{frame}
\frametitle{Dynamic complexity: motivation}
\begin{block}{Dynamic complexity}
How fast can we update the output when the input is slightly changed?
\end{block}
Why we are interested in dynamic complexity of the CFL-reachability problem?

\begin{itemize}
\item CFL-reachability is a dynamic problem
\item When code changes, how fast can we update the result of the CFL-reachability based code analysis?
\item New algorithmic approaches: expressing an updating in first-order logic, usage of the regular Datalog for updating, ...
\item Recently we have obtained reduction from the CFL-reachability problem to incremental transitive closure
\end{itemize}


\end{frame}

%-------------------------------------------------------------------------------------------------------------
\begin{frame}
\frametitle{Grants and students}
\begin{itemize}
\item Participating in research project ``Logical and algebraic methods in formal language theory'' (under supervision of Alexander Okhotin) funded by the Russian Scientific Foundation
\item Students: 
\begin{itemize}
\item Alexandra Olemskaya (HSE, now SPBSU) --- had her bachelor degree in HSE with thesis ``On some special cases of the CFL-reachability problem''
\item Alexandra Istomina (SPBSU)  --- working on her master thesis ``Fine-grained reductions around the CFL-reachability problem''
\end{itemize}
\item Teaching: ``Dynamic graph algorithms'' block (Spring 2021) in the ``Graph theory'' course (SPSU)
\end{itemize}
\end{frame}
%-------------------------------------------------------------------------------------------------------------
\begin{frame}
\frametitle{Current \& Future work}
\begin{itemize}
\item Dynamic complexity of the CFL-reachability in static and dynamic setting
\item Fine-grained complexity: add to the current reduction's zoo some new reductions
\item Generalization of obtained results on the CFL-reachability to more complicated language-constraint reachability problems (MCFG, Interleaved Dyck, ...)
\end{itemize}

\end{frame}
\end{document}

